\chapter*{KẾT LUẬN}
\addcontentsline{toc}{chapter}{KẾT LUẬN}

Qua đợt thực tập, em đã tìm hiểu được tổng quan về \CompanyName{} -- từ quá
trình hình thành, các lĩnh vực hoạt động cho đến cơ cấu tổ chức. Phần em tập
trung nhiều nhất là hoạt động cung ứng linh kiện điện tử và thiết bị tự động
hóa. Qua đó, em hiểu rõ hơn vai trò của việc đọc đúng mã hàng, đối chiếu
datasheet và giữ thông tin kỹ thuật nhất quán giữa các bộ phận liên quan.

Kết quả cụ thể của quá trình tìm hiểu là em đã phân loại bốn nhóm sản phẩm trong
danh mục TECHPAL, lập bảng kiểm trước khi lựa chọn mã hàng và tra cứu bốn mã đại
diện gồm MAX487CSA+T, BMXAMI0410, 6ES7153-1AA03-0XB0 và CR1000X. Bài toán giám
sát nhiệt độ tủ kỹ thuật giúp em thử đi từ yêu cầu ban đầu đến sơ đồ khối, sau
đó xem từng loại thiết bị có thể đảm nhận vai trò nào. Phương án này mới dừng ở
mức khảo sát trên giấy, chưa có nguyên mẫu và số liệu đo thực tế, nhưng giúp em
nhìn rõ hơn sự khác nhau giữa một IC giao tiếp, mô-đun PLC và bộ ghi dữ liệu
chuyên dụng.

Bài học lớn nhất em rút ra là không thể chọn linh kiện chỉ dựa vào tên gọi hoặc
mô tả ngắn trên trang bán hàng. Ngoài thông số điện, còn phải kiểm tra kiểu vỏ,
đầu nối, khả năng tương thích, điều kiện môi trường và trạng thái vòng đời của
sản phẩm. Chỉ một hậu tố bị bỏ sót cũng có thể dẫn đến sai quy cách đóng gói hoặc
sai biến thể cần đặt mua.

Do thời gian và điều kiện tiếp cận còn hạn chế, báo cáo chưa đi sâu vào quy trình
nội bộ, kiểm tra hàng hóa hoặc thử nghiệm một hệ thống hoàn chỉnh. Nếu có điều
kiện tiếp tục, em muốn hoàn thiện bảng so sánh cho các mã đã chọn và làm một mô
hình đo nhiệt độ nhỏ để kiểm tra lại phương án trên thực tế. Những kinh nghiệm
từ đợt thực tập sẽ là cơ sở để em học tốt hơn các nội dung về thiết kế mạch, hệ
thống nhúng và tự động hóa trong thời gian tới.

